\section{Physical reflection and the reference color junction}\label{sec:reflection}

\subsection{Ordinary time reflection}

Reflect a time direction tangent to the defect:
\begin{equation}
 \vartheta x=Rx,\qquad R=\diag(-1,1,1,1),\qquad
 b=0\in\{x^0=0\}\cap\{x^3=0\}.
\end{equation}
The positive half-space is $x^0>0$. In the bosonic adjoint convention,
$A$ transforms as a one-form, all six Hermitian scalars are even under
geometric time reflection, and conjugation exchanges $q_m$ with
$\bar q_m$. Fermionic reflection and an interacting regulator are
additional requirements, not consequences of these bosonic rules.

Let $C:a\to b$ lie in the positive half-space except at $b$, with $a$
on the defect. Use $M$ from \eqref{eq:M-map}. Pullback of the fields
first gives $U_{MR}(\vartheta C)$; applying \eqref{eq:dagger} then gives
\begin{equation}\label{eq:reflection-map}
 \Theta_0U_M(C)=U_N(\vartheta C^{-1}),\qquad N=-MR.
\end{equation}
Explicitly,
\begin{equation}
 Ne_0=-e_7,\quad Ne_1=-e_8,\quad Ne_2=-e_9,\quad Ne_3=-e_6.
\end{equation}
The time component retains its scalar map and the spatial components
reverse it. This ordering calculation agrees with the distinction
between geometric reflection and scalar-coupled adjoint transport
in \cite[eqs.~(4)--(7)]{DP}.

\subsection{The actual paired operator and the conditional Gram statement}

Define a covariant color vector at the reference by
\begin{equation}
 F_i^\alpha=[U_M(C_i)]^\alpha{}_{\beta}q_2^\beta(a_i).
\end{equation}
Its reflected row is
$\bar q_{2,\gamma}(\vartheta a_i)
[U_N(\vartheta C_i^{-1})]^\gamma{}_{\alpha}$. The proposed kernel is
\begin{equation}\label{eq:network}
 K_{ij}=\vev{
 \bar q_{2,\gamma}(\vartheta a_i)
 [U_N(\vartheta C_i^{-1})]^\gamma{}_{\alpha}
 [U_M(C_j)]^\alpha{}_{\beta}q_2^\beta(a_j)}.
\end{equation}
There is an identity color contraction at $b$, not another defect
field there. Under a gauge transformation, $F_i\mapsto g(b)F_i$ and
its reflected row transforms by $g(b)^{-1}$, so the paired observable
is gauge invariant. The individual vector is not a gauge-invariant
state without a specified reference charge sector.

\begin{assumption}\label{ass:OS}
A regulator and reference color sector admit Osterwalder--Schrader
positivity for these half-space functionals. The regulated endpoint
and junction operators belong to that domain, and their limiting
renormalization preserves the inequality whenever a limit is claimed.
\end{assumption}

\begin{proposition}[Conditional positivity]\label{prop:conditional-gram}
Under the regulated part of Assumption~\ref{ass:OS}, the matrix
in \eqref{eq:network} is positive semidefinite for every finite
collection of admissible contours and endpoints.
\end{proposition}
\begin{proof}
Apply the assumed OS inequality to $\sum_i c_iF_i$, including its
reference color fiber, to obtain $\sum_{ij}\bar c_iK_{ij}c_j\geq0$.
For a concrete bosonic sufficient model, suppose the two half-integrals
factor over boundary data $B$ with positive measure $\nu$. They give
vector wavefunctions $\psi_i^\alpha(B)$ and
\begin{equation}\label{eq:factorized-gram}
 K_{ij}=\int\sum_\alpha\overline{\psi_i^\alpha(B)}
                    \psi_j^\alpha(B)\dd\nu(B),\qquad\nu\geq0.
\end{equation}
The claimed inequality is then the integral of a squared norm.
Boundary gauge transformations act unitarily on the displayed fiber.
\end{proof}

This conditional statement locates the missing physical input. A
positive bosonic factorization is not asserted for the full fermionic
$\mathcal N=4$ defect measure. Gauge invariance of
\eqref{eq:network} alone proves no reflection inequality. Gauss law,
the reference charge sector, the regulator, and junction counterterms
must all be addressed in an interacting realization.

\subsection{An unconditional free kernel}

At leading free endpoint order $U=I$. Strip the positive color and
coupling constants and write $a_i=(\tau_i,\mathbf z_i,0)$ with
$\tau_i>0$, $\mathbf z_i\in\R^2$. Then
\begin{equation}\label{eq:free-OS}
 K^{(0)}(a_i,a_j)=
 \frac1{\sqrt{(\tau_i+\tau_j)^2+|\mathbf z_i-\mathbf z_j|^2}}.
\end{equation}
Its representation
\begin{equation}
 K^{(0)}(a_i,a_j)=4\pi\int_{\R^2}\frac{\mathrm d^2k}{(2\pi)^2}
 \frac{e^{-|k|(\tau_i+\tau_j)}
 e^{\ii k\cdot(\mathbf z_j-\mathbf z_i)}}{2|k|}
\end{equation}
is an integral of products $\overline{\psi_i(k)}\psi_j(k)$ with
positive measure. This proves free positivity independently of
Assumption~\ref{ass:OS}. On the time ray $a_r=(r,0,0,0)$ it reduces to
\begin{equation}\label{eq:opposite-ray}
 K^{(0)}(r,s)=\frac1{r+s}
 =\int_0^\infty e^{-\lambda r}e^{-\lambda s}\dd\lambda,
 \qquad
 \sqrt{rs}K^{(0)}(r,s)=G_o(\log r-\log s).
\end{equation}
The diagonal is finite for $r>0$. The kernel is Hankel in $r,s$,
and translation invariant after the radial weights and logarithmic
coordinates. It is the opposite-ray kernel of \eqref{eq:rays}, not
the full even archimedean tower. The full-sphere profile in
Section~\ref{sec:free} cannot be used as the same half-space state
without changing its support and recomputing its weights.

\subsection{The two branches preserve different constant charges}

Let $\widetilde K_\mu=\ii\Gamma_{Me_\mu}\Gamma_\mu$ for
$\mu=0,1,2$, and $K_n=\ii\Gamma_6\Gamma_3$. With $\chi=D=+1$,
the two branches require the following eigenvalues:
\begin{center}
\begin{tabular}{lcc}
\toprule
Involution & Ket with $M$ & Reflected bra with $N$\\
\midrule
$\widetilde K_0$ & $-1$ & $-1$\\
$\widetilde K_1,\widetilde K_2$ & $-1$ & $+1$\\
$K_n$ & $-1$ & $+1$\\
\bottomrule
\end{tabular}
\end{center}
Each all-direction family has one complex compatible charge; in
the $(x^0,x^3)$-plane each has two. Their common subspace is zero
when time and normal tangents are independently required, because
the normal conditions are opposite. The normal projectors also
select opposite H weights: the ket uses $q_2$ and the physical bra
uses $\bar q_2$ with the conjugate-weight charge. The identity tensor
at $b$ cannot compensate a nonzero variation on a branch.

As a bulk control, a straight time-directed contour needs only
$\widetilde K_0=-1$. That condition is common to the maps and gives
a four-dimensional complex chiral defect solution space. It is a
bulk control, not a statement that a constant charge annihilates
the physical adjoint pair of endpoint scalars.

\subsection{Internal rotations do not automatically preserve positivity}

The proper rotations by $\pi$ about H axis 4 and V axis 7 are
\begin{equation}
 S_H=\diag(1,-1,-1),\qquad S_V=\diag(1,-1,-1).
\end{equation}
For $S=S_H\oplus S_V$, one checks $SN=M$. Thus an internally rotated
reflection aligns the two bulk scalar maps. But $S_H$ acts on the
endpoint doublet by $\pm\ii\sigma_1$, up to basis convention. It
takes $\bar q_2$ to a phase times $\bar q_1$, restoring the zero
diagonal contraction of Proposition~\ref{prop:endpoint-zero}.

On the full free endpoint doublet the inserted metric is not positive.
Without a phase adjustment its eigenvalues are $\pm\ii$; with a phase
making it Hermitian they are $+1,-1$. For instance
\begin{equation}
 (1,-1)\sigma_1\binom{1}{-1}/2=-1.
\end{equation}
No overall phase makes both eigenvalues positive. The identity metric
from ordinary adjoint pairing is the control. More generally, if a
unitary $S$ on a positive Hilbert space satisfies
$\ip\psi{S\psi}\geq0$ for every $\psi$, then $S$ is positive
self-adjoint; unitarity forces its spectrum to be $\{1\}$ and $S=I$.
Projection to one rotation eigenspace is a different operation and
here mixes the fixed H endpoint weights.

This is a counterexample for the displayed doublet state space.
It does not contradict the special large-rectangle rotated-reflection
conjecture distinguished from ordinary reflection in
\cite[eqs.~(11)--(16)]{DP}.

\subsection{A free test of an odd-H normal reflection}

A different proposal in the parent investigation reflected across
the defect and assigned odd parity to the H scalars. For the canonical
free scalar algebra, a real mode of frequency $\nu>0$ at distance
$a>0$ from a reflection plane has ordinary norm
$e^{-2\nu a}/(2\nu)>0$; assigning it odd reflection parity changes
the sign. The failure is also visible in a gauge-invariant observable.
For a non-Abelian group with $\tr(T_aT_b)=\delta_{ab}/2$, take the
Hermitian Wick polynomial
\begin{equation}
 \mathcal O=\ii\tr\bigl(H[V_1,V_2]\bigr),
\end{equation}
where $H$ is one H scalar and $V_1,V_2$ are distinct V scalars.
Writing $[T_b,T_c]=\ii f_{bca}T_a$, its coefficient is
$-f_{bca}/2$. Distinct species have only their matching free
contractions, hence at separated reflected points
\begin{equation}
 \vev{\Theta_{\rm even}\mathcal O(x)\,\mathcal O(x)}
 =\tfrac14\sum_{abc}f_{abc}^2D_0(\vartheta x-x)^3>0,
\end{equation}
with $D_0$ the positive scalar covariance. Odd H and even V parity
give the negative of this value. This excludes that reflection on
the unrestricted canonical free bulk observable algebra. It does
not classify all interacting normal-reflection actions or restricted
observable sectors. The even assignment supplies the positive control.
